\begin{exo}{}
	Dans le plan muni du repère $\vOIJ$ orthonormé, on considère les trois points $A$, $B$ et $C$ de coordonnées : \cmath{A\coordl{-3}{-2} \quad ; \quad B\coordl{2}{0} \quad ; \quad C\coordl{-1}{3}}

	\begin{enumerate}
		\item
		      \begin{enumerate}
			      \item Déterminer les coordonnées du vecteur $\vec{AB}$.
			      \item En déduire que : $\norm{\vec{AB}} = \sqrt{29}$.
		      \end{enumerate}
		\item Déterminer la norme $\norm{\vec{AC}}$.
		\item Placer les points $A$, $B$ et $C$ dans le repère ci-dessous :

		      \begin{center}
			      \begin{tikzpicture}[scale=0.80, >=Stealth, font=\footnotesize]
				      \draw[gray!30, very thin, step=0.2] (-4.5,-2.5) grid (4.5,3.5);
				      % Quadrillage principal
				      \draw[gray!60, thin, step=1] (-4.5,-2.5) grid (4.5,3.5);

				      % Axes
				      \draw[thick, ->] (-4.5,0) -- (4.5,0);
				      \draw[thick, ->] (0,-2.5) -- (0,3.5);

				      % Labellisation des axes
				      \foreach \x in {-4,-3,-2,-1,1,2,3,4} { \draw (\x,0.1) -- (\x,-0.1) node[below] {$\x$}; }
				      \foreach \y in {-2,-1,1,2,3} {
						      \draw (0.1,\y) -- (-0.1,\y) node[left] {$\y$};
					      }

				      % Unités de base O, I, J
				      \node[below left, font=\footnotesize] at (0,0) {$O$};
			      \end{tikzpicture}
		      \end{center}

		\item Justifier que $ABC$ est isocèle en $A$.
	\end{enumerate}
\end{exo}
